\documentclass[runningheads]{llncs}
\usepackage[T1]{fontenc}
\usepackage{graphicx}
\usepackage{tabularx}
\usepackage{xurl}
\usepackage{booktabs}
\usepackage{array}
\usepackage{threeparttable}
\usepackage{ragged2e}
\setlength{\textfloatsep}{8pt plus 2pt minus 2pt}
\setlength{\intextsep}{8pt plus 2pt minus 2pt}
\setlength{\floatsep}{6pt plus 2pt minus 2pt}
\setlength{\emergencystretch}{2em}


\begin{document}
\title{Integrating Book Data at Edition Level: Goodreads, DBpedia and Amazon}
\titlerunning{Book Data Integration}
\authorrunning{Bolduan et al.}
\author{Lucia Bolduan \orcidID{2261625} \and Safa Shaghely \orcidID{2480414} \and Selma Nezihoglu \orcidID{2390917} \and Minji Oh \orcidID{2389125} \and Kejsi Limaci \orcidID{2389812} \and Shiqi Xu \orcidID{1863435}}
\institute{University of Mannheim}
\maketitle

\begin{abstract}
% TODO
\end{abstract}

\section{Introduction and Data}
% TODO: use case, entity = edition, the three sources, integrated schema (short version of the proposal)

\section{Phase 1: Schema Mapping and Data Translation}

\subsection{Target Schema}
The integrated schema is defined as an XML Schema (\texttt{target.xsd}) and
serves as the common target for all three sources. Its root element
\texttt{books} contains a sequence of \texttt{book} elements, each describing
exactly one book edition, in line with our entity definition. Every book
carries a required \texttt{id} attribute that is prefixed with the name of its
source (e.g.\ \texttt{goodreads\_id\_1}). These IDs keep each record traceable
to its origin and serve as keys in the correspondences produced during
identity resolution.

The schema contains 14 attributes. Only \texttt{title} is mandatory; all other
elements are optional (\texttt{minOccurs="0"}). We chose this deliberately: no
source covers all attributes, and a stricter schema would either reject valid
records or force us to insert artificial default values that would later
distort data fusion. Missing values are therefore simply omitted.

We assigned explicit data types so that later steps can compare values without
additional parsing: \texttt{page\_count} and \texttt{ratings\_count} are
\texttt{xs:integer}, \texttt{rating} and \texttt{price} are
\texttt{xs:decimal}, and \texttt{publication\_date} is \texttt{xs:date}.
Multi-valued attributes are modelled as nested lists instead of delimited
strings: \texttt{authors}, \texttt{genres} and \texttt{awards} contain one
child element per value. This enables list-based comparators during identity
resolution and union-based fusion of genres. Compared with our proposal,
\texttt{format} and \texttt{edition} were added because they help to
distinguish editions of the same work, and the full publication date is kept
instead of only the year.

The \texttt{series} attribute remains in the target schema but is filled
only from Goodreads, the only source with a dedicated field (23,469
records, 44.7\%). In Amazon, 38.4\% of the titles contain brackets, but only
a small share (about 3\%) clearly looks like a numbered series; most
brackets describe the edition instead, such as ``Spanish Edition'' or
``Random House Large Print''. In DBpedia, 21.3\% of the titles contain
brackets, almost exclusively Wikipedia disambiguations such as ``(novel)''
(6,446) or ``(book)'' (1,857). Since the series cannot be extracted
reliably for Amazon and DBpedia, we do not use it for identity resolution
or data fusion.

\subsection{Mappings}
We implemented one Altova MapForce mapping per source. Goodreads and Amazon
are read as CSV files, DBpedia as the XML file produced by our SPARQL download
script. Each mapping writes an XML file that is valid against
\texttt{target.xsd}. Goodreads provides the richest attribute set, including
series, awards, rating, ratings count and price, but stores several values in
single strings. Amazon provides format and edition as separate fields and
lists authors and genres as strings that need to be split. DBpedia
contributes title, authors, publisher, release date, pages, language and
genres. Each record receives a source-prefixed ID generated with
\texttt{auto-number} and \texttt{concat}. Because these IDs depend on the
order of the output, the XML files must not be regenerated without re-running
identity resolution; otherwise the stored correspondences point to the wrong
records.

\subsection{User-Defined Functions}
Many cleaning steps are needed in more than one mapping. Instead of rebuilding
the same chains of built-in functions, we implemented reusable user-defined
functions (UDFs), which keep the mappings readable and ensure that the same
rule is applied everywhere (Table~\ref{tab:udf}).

\begin{table}[htbp]
\centering
\caption{User-defined functions in the MapForce mappings.}
\label{tab:udf}
\footnotesize
\begin{tabularx}{\textwidth}{@{}l l >{\RaggedRight\arraybackslash}X@{}}
\toprule
\textbf{Function} & \textbf{Used for} & \textbf{Purpose} \\
\midrule
\texttt{Clean\_Authors} & GR author & splits the author string into individual author elements \\
\texttt{Clean\_List} & GR genres, awards; AM authors & splits multi-valued strings and drops empty entries \\
\texttt{Parse\_Date} & GR publishDate & detects four date formats and converts them to \texttt{xs:date} \\
\texttt{Page\_Count} & pages & extracts the number from values such as ``1 page'' \\
\texttt{Safe\_Number} & GR price & converts only valid numbers; other values stay empty \\
\bottomrule
\end{tabularx}
\end{table}

\subsection{Date Formats}
The sources differ strongly in date representation and precision
(Table~\ref{tab:dates}).

\begin{table}[htbp]
\centering
\caption{Precision of the publication date per source.}
\label{tab:dates}
\footnotesize
\begin{tabular}{lrrr}
\toprule
 & \textbf{Goodreads} & \textbf{DBpedia} & \textbf{Amazon} \\
\midrule
Full date               & 84.1\% & 22.2\% & 95.1\% \\
Year only               &  8.3\% & 58.0\% &  0.0\% \\
Other (e.g.\ month--year) & 5.9\% & 10.5\% & 0.0\% \\
Empty                   &  1.7\% &  9.2\% &  4.9\% \\
\bottomrule
\end{tabular}
\end{table}

Goodreads uses four formats: US dates with a two-digit year (04/27/10), with a
four-digit year (04/27/2010), written dates with an ordinal suffix (April 27th
2010) and bare years (2003). Amazon uses written dates such as January 1,
2019. DBpedia combines several infobox properties, so 58\% of its values
contain only a year, followed by ISO dates (1982-04-15) and month--year values
(July 1941).

\texttt{Parse\_Date} first identifies the Goodreads format with regular
expressions (\texttt{match-pattern}) and then applies a dedicated conversion
in a cascade of \texttt{if-else} components. Two-digit years are completed
with the century (above 26 $\rightarrow$ 19xx, otherwise 20xx), four-digit US
dates are parsed with the pattern \texttt{[M01]/[D01]/[Y0001]}, ordinal
suffixes are removed with \texttt{substring-before}, \texttt{substring-after}
and \texttt{string-length} before parsing with \texttt{[MNn] [D1] [Y0001]},
and bare years are completed to 1~January. Values matching none of the
patterns are left empty, so no invalid dates are produced. For Amazon, the
date is matched with a regular expression and converted with
\texttt{parse-date}. % TODO: DBpedia

Completing year-only values to 1~January makes these dates look more precise
than they are. In Amazon, 17.6\% of the dates fall on 1~January, which
suggests that the source itself fills unknown days with this default.

\subsection{Numeric Values and Further Problems}
A single invalid value in an \texttt{xs:integer} or \texttt{xs:decimal} field
makes MapForce abort the transformation. We therefore read numeric fields as
strings and convert them with UDFs. \texttt{Safe\_Number} converts the
Goodreads price only if it is a valid number; malformed values such as
1.189.88 (12 records) and 14,365 empty values are omitted, leaving 38,101
valid prices. \texttt{Page\_Count} extracts the number from values such as
``1 page'' (23 records), so that 50,131 Goodreads records carry a page count.
Splitting Amazon genres produced empty tokens, which we removed with a
filter. Finally, some Goodreads award strings contain escaped HTML entities
(\texttt{\&amp;}) that already exist in the source CSV; since they do not
affect matching, we left them unchanged.

% =====================================================================
\section{Phase 2: Identity Resolution}
% Authors: Shiqi Xu (Gold Standard, Blocking, Learning-Based Matching, Global Matching),
% Kejsi Limaci (Rule-Based Matching, Results and Error Analysis)
% All numbers refer to commit 4daa279 (identity_resolution.ipynb). Do not re-run
% the matching after this commit, otherwise the numbers and the fusion evaluation
% set no longer match.

We match each of the three dataset pairs separately with PyDI. The pipeline
consists of a gold standard (Sect.~\ref{sec:gs}), blocking
(Sect.~\ref{sec:blocking}), rule-based and learning-based matchers
(Sect.~\ref{sec:rules}, \ref{sec:ml}), global one-to-one matching with a
consistency check across the three pairs (Sect.~\ref{sec:global}) and an
evaluation on held-out test pairs (Sect.~\ref{sec:irresults}).

% ---------------------------------------------------------------- Shiqi
\subsection{Gold Standard}\label{sec:gs}
Since none of the sources links its records to the others, we derive the labels
from the ISBN, which is not used by any matcher (weak supervision). Two records
are labelled as a match if they share an ISBN-13. Following the lecture, the
gold standard contains matches and non-matches of different difficulty
(Table~\ref{tab:gs}): \emph{random matches}; \emph{corner-case matches}, whose
titles differ strongly (word Jaccard $<0.5$, e.g.\ a subtitle or a series name
in only one source); \emph{corner-case non-matches}, i.e.\ different editions of
the same work (same normalised title and author, different ISBN); and
\emph{random non-matches}. The random non-matches are drawn from pairs by the
same author instead of from all pairs, because completely random pairs are
trivially different and would make the evaluation too optimistic (selection
bias).

\begin{table}[htbp]
\centering
\caption{Gold standard per dataset pair (after the manual check).}
\label{tab:gs}
\footnotesize
\begin{tabular}{lrrrrrr}
\toprule
 & \multicolumn{2}{c}{\textbf{Matches}} & \multicolumn{2}{c}{\textbf{Non-matches}} & & \\
\cmidrule(lr){2-3}\cmidrule(lr){4-5}
\textbf{Pair} & random & corner & corner & random & \textbf{Train} & \textbf{Test} \\
\midrule
Goodreads--Amazon  & 300 & 65  & 375 & 450 & 833 & 357 \\
Goodreads--DBpedia & 245 & 142 & 371 & 450 & 845 & 363 \\
Amazon--DBpedia    & 258 & 59  & 371 & 450 & 796 & 342 \\
\bottomrule
\end{tabular}
\end{table}

An ISBN is not a perfect label, so we checked 1,428 sampled pairs. Of 1,112
ISBN-based matches, 239 (21\%) did not describe the same edition, while only 8
of 316 non-matches were wrong. The main reason is DBpedia: its infobox describes
the first edition, but the ISBN often belongs to a later edition, so a shared
ISBN does not mean that the two records describe the same edition (108 of the
249 removed pairs). Other errors were different books with a shared ISBN (e.g.\
a box set and a single volume) and translations. The 249 wrong pairs are listed
in \texttt{excluded\_pairs.csv} and removed. Each gold standard is split
70/30 into training and test set; 25\% of the training pairs serve as
validation set for all choices (model, hyperparameters, threshold), so the test
set is used only once.

\subsection{Blocking}\label{sec:blocking}
Without blocking, each pair would require up to $3.1\cdot10^9$ comparisons.
We combine three exact-key blockers: the surname of the first author, the
cleaned title, and the cleaned title before a colon, matched against the full
title of the other source (``Mockingjay: The Final Book'' vs.\
``Mockingjay''). A pair becomes a candidate if it shares at least one key.
Table~\ref{tab:blocking} shows that the title alone loses up to half of the
matches, because titles differ in subtitles, series names and DBpedia
disambiguations, while the union keeps 96--99\% of them with a reduction ratio
above 99.9\%. The remaining loss consists mostly of author names written
differently (initials, transliterations) together with a changed title. These
matches count as false negatives in all results below.

\begin{table}[htbp]
\centering
\caption{Blocking: pair completeness (PC) and number of candidates.}
\label{tab:blocking}
\footnotesize
\begin{tabular}{lrrr}
\toprule
\textbf{Pair} & \textbf{PC title} & \textbf{PC union} & \textbf{Candidates (union)} \\
\midrule
Goodreads--Amazon  & 72.4\% & 99.3\% & 1,315,284 \\
Goodreads--DBpedia & 53.4\% & 97.8\% & 1,071,227 \\
Amazon--DBpedia    & 65.4\% & 96.1\% & 1,212,964 \\
\bottomrule
\end{tabular}
\end{table}

% ---------------------------------------------------------------- Kejsi
\subsection{Rule-Based Matching}\label{sec:rules}
% TODO Kejsi (about half a page):
% - baseline rule (which comparators, threshold chosen on validation)
% - tuned linear weighted rule: weight grid, threshold, publisher_norm decision
% - non-linear rule: gate idea and variants
% - one sentence why rules struggle on the DBpedia pairs (test F1 in Table ref{tab:irresults})

% ---------------------------------------------------------------- Shiqi
\subsection{Learning-Based Matching}\label{sec:ml}
Each candidate pair is described by 16 similarity features: Jaro-Winkler and
Levenshtein on the cleaned title, word Jaccard on the full title, Jaro-Winkler
on the title before the colon and on the first author and Jaccard on the author lists,
Jaccard on the cleaned and on the normalised publisher, the absolute year
difference (up to 10 years) and the relative page-count difference (up to
50\%), and Jaro-Winkler on format and language. Since a missing value and a
dissimilar value both yield a low similarity, four indicator features state
whether year, page count, publisher and format are present in both records.
We trained Logistic Regression, Random Forest and Gradient Boosting, tuned their
hyperparameters by grid search with 5-fold stratified cross-validation on the
training set (F1) and selected the model with the highest F1 on the validation
set: Gradient Boosting for Goodreads--Amazon (validation F1 0.90) and Random
Forest for the two DBpedia pairs (0.78 and 0.77). The tree ensembles outperform
Logistic Regression mainly because the features interact: a high title
similarity indicates a match only if year and page count also agree, otherwise
it describes another edition of the same work.

\subsection{Global Matching and Cluster Consistency}\label{sec:global}
The lecture recommends local matching for duplicate detection within one source
and global matching for sources without duplicates. Our sources are nearly
duplicate-free at edition level: only 0.2\% of the Goodreads, 2.8\% of the
Amazon and 0.8\% of the DBpedia records share an ISBN with another record of
the same source. We therefore apply greedy one-to-one matching to the match
probabilities; only the threshold (0.5--0.9) is tuned on the validation set.

The gold standard cannot show how many wrong correspondences are produced
outside the labelled pairs. As a plausibility check, we estimate the number of
true matches from the shared ISBNs divided by the share of records with an ISBN
in both sources (2,642, 6,590 and 2,854). Without the one-to-one constraint the
matchers return 11,251, 17,465 and 5,801 correspondences, four times the
estimate for Goodreads--Amazon, although the test F1 is similar or even higher
for the DBpedia pairs (Table~\ref{tab:irresults}). With the constraint the
numbers (4,146, 7,369 and 2,362) are much closer to the estimate.

One-to-one matching holds within each pair, but chaining the three pairs still
produced clusters with up to ten records (4.5\% of the clusters have four or
more), i.e.\ two records of the same source were connected through a third
record. Since each source lists an edition at most once, such a cluster contains
at least one wrong correspondence. We therefore process all correspondences of
the three pairs in descending order of their score and keep a correspondence
only if it does not put two records of the same source into one cluster. This
removes 3--6\% of the correspondences and costs about one point F1
(Table~\ref{tab:irresults}), but guarantees that data fusion never merges two
editions of the same source. The final 13,340 correspondences form 9,960
clusters of size two and 1,433 of size three.

% ---------------------------------------------------------------- Kejsi
\subsection{Results and Error Analysis}\label{sec:irresults}

\begin{table}[htbp]
\centering
\caption{Test results (F1; precision/recall for the final configuration).}
\label{tab:irresults}
\footnotesize\setlength{\tabcolsep}{4pt}
\begin{tabular}{lrrrrc}
\toprule
 & \multicolumn{2}{c}{\textbf{Rules}} & \multicolumn{3}{c}{\textbf{ML}} \\
\cmidrule(lr){2-3}\cmidrule(lr){4-6}
\textbf{Pair} & linear & non-lin. & no 1:1 & 1:1 & 1:1 + cleanup: P / R / \textbf{F1} \\
\midrule
Goodreads--Amazon  & 0.765 & 0.774 & 0.851 & 0.851 & 0.961 / 0.753 / \textbf{0.844} \\
Goodreads--DBpedia & 0.688 & 0.669 & 0.760 & 0.737 & 0.758 / 0.694 / \textbf{0.725} \\
Amazon--DBpedia    & 0.535 & 0.547 & 0.804 & 0.750 & 0.843 / 0.663 / \textbf{0.742} \\
\bottomrule
\end{tabular}
\end{table}

% TODO Kejsi (about three quarters of a page):
% - compare rules vs. ML (Table ref{tab:irresults}); why the DBpedia pairs are hardest
% - typical errors with one example each, e.g.
%     * DBpedia first edition vs. later edition (label noise, see Sect. ref{sec:gs})
%     * same title, different book/author (e.g. "Fallen" by K. George vs. K. Slaughter)
%     * volumes of a series / omnibus editions
%     * missing or damaged values (e.g. dates lost in the Goodreads source)
% - recall loss: blocking (Table ref{tab:blocking}) + one-to-one + cleanup
% - future work: e.g. more corner cases in training, edition-aware DBpedia handling

% =====================================================================
\section{Phase 3: Data Fusion}
% TODO

\section{Conclusion}
% TODO

% \bibliographystyle{splncs04}
% \bibliography{references}
\end{document}
